\documentclass[11pt,a4paper]{article}
\usepackage[utf8]{inputenc}
\usepackage[T1]{fontenc}
\IfFileExists{romanian.ldf}{\usepackage[romanian]{babel}}{}
\usepackage{lmodern}
\usepackage[margin=2.2cm]{geometry}
\usepackage{xcolor}
\usepackage{listings}
\usepackage{booktabs}
\usepackage{enumitem}
\usepackage{fancyhdr}
\usepackage[most]{tcolorbox}
\usepackage[hidelinks]{hyperref}

\definecolor{codebg}{RGB}{247,247,250}
\definecolor{kw}{RGB}{0,70,150}
\definecolor{cm}{RGB}{90,120,90}
\definecolor{st}{RGB}{160,60,30}

\lstset{
  basicstyle=\ttfamily\small,
  backgroundcolor=\color{codebg},
  keywordstyle=\color{kw}\bfseries,
  commentstyle=\color{cm}\itshape,
  stringstyle=\color{st},
  numbers=left, numberstyle=\tiny\color{gray},
  frame=single, rulecolor=\color{gray!40},
  breaklines=true, columns=fullflexible, keepspaces=true,
  showstringspaces=false, tabsize=4,
  literate={ă}{{\u{a}}}1 {â}{{\^{a}}}1 {î}{{\^{\i}}}1 {ș}{{\c{s}}}1 {ț}{{\c{t}}}1
}

\newtcolorbox{terminal}[1][]{colback=black!85,colframe=black!60,coltext=white,
  fontupper=\ttfamily\small,boxrule=0.4pt,arc=2pt,left=4pt,right=4pt,top=3pt,bottom=3pt,title={#1},
  fonttitle=\sffamily\small\bfseries,colbacktitle=black!60}
\newtcolorbox{answer}{colback=blue!3,colframe=blue!40!black,boxrule=0.4pt,arc=2pt,
  left=6pt,right=6pt}

\pagestyle{fancy}
\fancyhf{}
\lhead{Sisteme de Operare -- FCIM, UTM}
\rhead{Ciur Alexandrina}
\cfoot{\thepage}

\begin{document}

\begin{titlepage}
\centering
{\large Ministerul Educației și Cercetării al Republicii Moldova\\
Universitatea Tehnică a Moldovei\\
Facultatea Calculatoare, Informatică și Microelectronică\par}
\vspace{5cm}
{\LARGE\bfseries RAPORT\par}
\vspace{0.6cm}
{\Large Lucrarea de laborator nr.~13\par}
\vspace{0.4cm}
{\large la disciplina \textit{Sisteme de Operare}\par}
\vspace{0.6cm}
{\Large\bfseries Semaphores in Your Own Kernel (xv6)\par}
\vfill
\begin{flushright}
\begin{tabular}{ll}
A efectuat: & Ciur Alexandrina \\
A verificat: & \rule{4.5cm}{0.4pt} \\
\end{tabular}
\end{flushright}
\vspace{2cm}
{\large Chișinău, 2026\par}
\end{titlepage}

\section*{Scopul lucrării}
În laboratoarele 11 și 12 semafoarele au fost \emph{folosite} din biblioteca C. Acum le \emph{construim} în kernelul xv6,
ca trei apeluri de sistem noi: \texttt{sem\_init}, \texttt{sem\_down} și \texttt{sem\_up} (numerele 24, 25, 26). Scopul e
să vedem că un semafor e doar un număr, un spinlock care îl protejează și mecanismul \texttt{sleep}/\texttt{wakeup} al
kernelului. Procesele care așteaptă dorm cu adevărat și nu consumă CPU.

\section*{Noțiuni teoretice}
\begin{itemize}[leftmargin=*]
  \item \textbf{Apel de sistem în xv6}: programul utilizator apelează un \emph{stub} generat de \texttt{usys.pl}, care
        pune numărul apelului în registrul \texttt{a7} și execută \texttt{ecall}. Kernelul caută numărul în tabelul
        \texttt{syscalls[]} din \texttt{syscall.c} și apelează funcția \texttt{sys\_*}.
  \item \textbf{Spinlock}: lock în care CPU-ul se învârte (busy-wait) până devine liber. În xv6, \texttt{acquire}
        dezactivează și întreruperile. Se ține doar foarte scurt și niciodată în timp ce procesul doarme.
  \item \textbf{sleep / wakeup}: un proces adoarme pe un \emph{canal} (o adresă oarecare), iar \texttt{wakeup(chan)}
        îi face RUNNABLE pe toți cei care dorm pe acel canal. Kernelul nostru adoarme în doi pași:
        \texttt{sleep\_prepare(chan)}, apoi \texttt{sleep()}. Dacă între ei vine un \texttt{wakeup},
        \texttt{sleep()} se întoarce imediat, deci trezirea nu se pierde.
\end{itemize}

%======================================================================
\section{Partea 1 -- Numerele și partea de utilizator}
Rețeta e aceeași ca la \texttt{trace} din laboratorul 10, aplicată de trei ori. Semafoarele sunt numerotate de la 0 la 7
și sunt comune tuturor proceselor. Numărul 23 rămâne rezervat pentru \texttt{trace}.

\begin{lstlisting}[language=C,title={\texttt{snippets/syscall.h}}]
// kernel/syscall.h
#define SYS_sync     22
#define SYS_sem_init 24   // new
#define SYS_sem_down 25   // new
#define SYS_sem_up   26   // new
\end{lstlisting}
\begin{lstlisting}[language=C,title={\texttt{snippets/user.h}}]
// user/user.h
int sync(void);
int sem_init(int, int);   // new: id, start value
int sem_down(int);        // new: take one, or sleep
int sem_up(int);          // new: give one back
\end{lstlisting}
\begin{lstlisting}[language=Perl,title={\texttt{snippets/usys.pl}}]
# user/usys.pl
entry("sync");
entry("sem_init");       # new
entry("sem_down");       # new
entry("sem_up");         # new
\end{lstlisting}
\begin{lstlisting}[language=make,title={\texttt{snippets/Makefile}}]
# Makefile - OBJS (two spaces in front)
  $K/sysproc.o \
  $K/sem.o \

# Makefile - UPROGS (a TAB in front)
	$U/_zombie\
	$U/_semtest\
	$U/_mix\
\end{lstlisting}

\subsection*{Programul de test \texttt{semtest}}
Testul 1: copilul trebuie să aștepte părintele (semaforul 0 pornește de la 0). Testul 2: ping și pong alternează, cu două
semafoare care își pasează rândul.
\begin{lstlisting}[language=C,title={\texttt{user/semtest.c}}]
// semtest.c: two tests for the kernel semaphores (Lab 13).
#include "kernel/types.h"
#include "user/user.h"

int
main(void)
{
  // Test 1: order. Semaphore 0 starts at 0, so the child must wait.
  sem_init(0, 0);
  if (fork() == 0) {
    sem_down(0);
    printf("child: I ran after the parent\n");
    exit(0);
  }
  pause(10);
  printf("parent: done, I let the child go\n");
  sem_up(0);
  wait(0);

  // Test 2: ping-pong. Semaphore 1 = ping's turn, semaphore 2 = pong's turn.
  sem_init(1, 1);
  sem_init(2, 0);
  if (fork() == 0) {
    for (int i = 1; i <= 3; i++) {
      sem_down(2);
      printf("  pong %d\n", i);
      sem_up(1);
    }
    exit(0);
  }
  for (int i = 1; i <= 3; i++) {
    sem_down(1);
    printf("ping %d\n", i);
    sem_up(2);
  }
  wait(0);
  printf("semtest: OK\n");
  exit(0);
}
\end{lstlisting}

\subsection*{Programul \texttt{mix}}
\begin{lstlisting}[language=C,title={\texttt{user/mix.c}}]
// mix.c: two processes print five lines each, one letter per write().
//   mix        no lock: the letters of the two processes can mix
//   mix lock   semaphore 3 used as a lock: every line stays whole
#include "kernel/types.h"
#include "user/user.h"

void
lines(char *s, int uselock)
{
  for (int k = 0; k < 5; k++) {
    if (uselock)
      sem_down(3);
    for (char *p = s; *p; p++)
      write(1, p, 1);
    if (uselock)
      sem_up(3);
  }
}

int
main(int argc, char *argv[])
{
  int uselock = (argc > 1);
  sem_init(3, 1);
  if (fork() == 0) {
    lines("bbbbbbbbbbbbbbbbbbbbbbbbbbbbbb\n", uselock);
    exit(0);
  }
  lines("AAAAAAAAAAAAAAAAAAAAAAAAAAAAAA\n", uselock);
  wait(0);
  exit(0);
}
\end{lstlisting}

\begin{terminal}[Check 1: apelurile ajung în kernel, care încă nu le cunoaște]
\begin{verbatim}
$ semtest
3 semtest: unknown sys call 24
4 semtest: unknown sys call 25
child: I ran after the parent
parent: done, I let the child go
...
\end{verbatim}
\end{terminal}
Linia \texttt{child:} apare \textbf{înaintea} liniei \texttt{parent:}. Fără semafoare, copilul nu așteaptă nimic și rulează
primul, pentru că părintele face \texttt{pause(10)}. Mesajele \texttt{unknown sys call} arată că stub-urile funcționează,
dar tabelul din kernel nu are încă intrările 24--26.

%======================================================================
\section{Partea 2 -- Partea de kernel}
Fișierul nou \texttt{kernel/sem.c}. Inima lui e \texttt{sys\_sem\_down}: cât timp valoarea e 0, procesul se înregistrează
pe adresa semaforului, eliberează spinlock-ul și adoarme. \texttt{sys\_sem\_up} crește valoarea cu 1 și îi trezește pe toți
cei care dorm pe acea adresă.

\begin{lstlisting}[language=C,title={\texttt{kernel/sem.c}}]
// sem.c: counting semaphores inside the xv6 kernel (Lab 13).
// A small table of NSEM semaphores, shared by every process.
// Each one is a number plus a spinlock that protects the number.

#include "types.h"
#include "riscv.h"
#include "defs.h"
#include "param.h"
#include "spinlock.h"
#include "proc.h"

#define NSEM 8

struct sem {
  struct spinlock lock;   // a zeroed spinlock is a free spinlock
  int value;              // how many may still pass
};

static struct sem sems[NSEM];

// sem_init(id, value): give semaphore id its start value.
uint64
sys_sem_init(void)
{
  int id, value;
  argint(0, &id);
  argint(1, &value);
  if (id < 0 || id >= NSEM || value < 0)
    return -1;
  acquire(&sems[id].lock);
  sems[id].value = value;
  release(&sems[id].lock);
  return 0;
}

// sem_down(id): take one unit; if there is none, sleep until there is.
uint64
sys_sem_down(void)
{
  int id;
  argint(0, &id);
  if (id < 0 || id >= NSEM)
    return -1;
  struct sem *s = &sems[id];

  acquire(&s->lock);
  while (s->value == 0) {
    sleep_prepare(s);       // I will sleep on the address of s
    release(&s->lock);      // never sleep holding a spinlock
    sleep();                // returns at once if sem_up came in between
    acquire(&s->lock);
    if (killed(myproc())) { // kill must still work on a waiting process
      release(&s->lock);
      return -1;
    }
  }
  s->value--;
  release(&s->lock);
  return 0;
}

// sem_up(id): give one unit back and wake whoever sleeps on s.
uint64
sys_sem_up(void)
{
  int id;
  argint(0, &id);
  if (id < 0 || id >= NSEM)
    return -1;
  struct sem *s = &sems[id];

  acquire(&s->lock);
  s->value++;
  wakeup(s);
  release(&s->lock);
  return 0;
}
\end{lstlisting}

Fișierul e adăugat în lista \texttt{OBJS} din Makefile (\texttt{\$K/sem.o}), iar funcțiile sunt legate în tabelul de
apeluri de sistem:
\begin{lstlisting}[language=C,title={\texttt{snippets/syscall.c}}]
// kernel/syscall.c
extern uint64 sys_sync(void);
extern uint64 sys_sem_init(void);   // new
extern uint64 sys_sem_down(void);   // new
extern uint64 sys_sem_up(void);     // new
...
  [SYS_sync]     = sys_sync,
  [SYS_sem_init] = sys_sem_init,    // new
  [SYS_sem_down] = sys_sem_down,    // new
  [SYS_sem_up]   = sys_sem_up,      // new
\end{lstlisting}

\begin{terminal}[Check 2: semtest]
\begin{verbatim}
$ semtest
parent: done, I let the child go
child: I ran after the parent
ping 1
  pong 1
ping 2
  pong 2
ping 3
  pong 3
semtest: OK
\end{verbatim}
\end{terminal}
Părintele așteaptă 10 tick-uri înainte să-l elibereze pe copil, dar copilul afișează al doilea: a dormit în
\texttt{sem\_down(0)}. Ping și pong alternează strict, de câte trei ori.

%======================================================================
\section{Partea 3 -- Semaforul ca lock: \texttt{mix}}
\texttt{mix} pornește două procese care afișează câte cinci linii, câte o literă per \texttt{write()}. \texttt{make qemu}
folosește 3 CPU-uri, deci procesele rulează cu adevărat în paralel.

\begin{terminal}[Fără lock]
\begin{verbatim}
$ mix
AAAAAAAbAbbAbAbbbbbAbAbAbbAbAbbbbbbbbbbbbbbbb
AbAbbbbbbbAbbAbAbAbAbAbbbbbbbbbbbbbbbAb
AAAAbbAA
bbbbbbbbbbbAbAAAAbbbAbAbAbAbAbAbAbAbAbAbAbAbAbA
AbAbAbAbAbAbAbAbAbbbAbbAb
\end{verbatim}
\end{terminal}

\begin{terminal}[Cu lock (semaforul 3, pornit de la 1)]
\begin{verbatim}
$ mix lock
AAAAAAAAAAAAAAAAAAAAAAAAAAAAAA
AAAAAAAAAAAAAAAAAAAAAAAAAAAAAA
AAAAAAAAAAAAAAAAAAAAAAAAAAAAAA
AAAAAAAAAAAAAAAAAAAAAAAAAAAAAA
AAAAAAAAAAAAAAAAAAAAAAAAAAAAAA
bbbbbbbbbbbbbbbbbbbbbbbbbbbbbb
bbbbbbbbbbbbbbbbbbbbbbbbbbbbbb
bbbbbbbbbbbbbbbbbbbbbbbbbbbbbb
bbbbbbbbbbbbbbbbbbbbbbbbbbbbbb
bbbbbbbbbbbbbbbbbbbbbbbbbbbbbb
\end{verbatim}
\end{terminal}

\subsection*{Observații}
\begin{answer}
\textbf{1. De ce se amestecă literele fără lock.} O linie costă \textbf{31 de apeluri} \texttt{write()} (30 de litere plus
\texttt{\textbackslash n}). Între oricare două \texttt{write}-uri, celălalt proces (care rulează în paralel pe alt CPU,
sau după o comutare de context) poate scrie și el. Fiecare \texttt{write()} e indivizibil doar pentru litera lui, nu pentru
toată linia.

\medskip
\textbf{2. Liniile care fac o linie secțiune critică.} \texttt{if (uselock) sem\_down(3);} înainte de bucla de
\texttt{write} și \texttt{if (uselock) sem\_up(3);} după ea. Semaforul 3 pornește de la 1, deci e un lock: un singur proces
poate fi între ele, iar celălalt doarme în kernel.

\medskip
\textbf{3. De ce \texttt{while (s->value == 0)}, nu \texttt{if}.} \texttt{wakeup(s)} îi trezește pe \emph{toți} cei care
dorm pe \texttt{s}, dar \texttt{sem\_up} adaugă o singură unitate. Primul care reia spinlock-ul ia unitatea, iar ceilalți
găsesc tot 0. Un al treilea proces poate face \texttt{sem\_down} între \texttt{wakeup} și momentul în care cel trezit reia
lock-ul. Și un \texttt{kill} trezește procesul fără ca valoarea să se fi schimbat. \texttt{while} reverifică condiția și
procesul se culcă la loc.

\medskip
\textbf{4. De ce \texttt{release} înainte de \texttt{sleep()}.} \texttt{sem\_up} are nevoie de același spinlock ca să
facă \texttt{value++} și \texttt{wakeup}. Dacă procesul ar dormi ținând lock-ul, \texttt{sem\_up} s-ar învârti la
nesfârșit în \texttt{acquire} și ar apărea un deadlock: cel care doarme așteaptă o trezire care nu poate veni. În plus,
\texttt{acquire} dezactivează întreruperile, iar scheduler-ul xv6 nu permite cedarea CPU-ului cu un spinlock în plus ținut
(\texttt{panic("sched locks")}). Ordinea \texttt{sleep\_prepare}, \texttt{release}, \texttt{sleep} garantează că un
\texttt{sem\_up} venit între \texttt{release} și \texttt{sleep()} nu se pierde.
\end{answer}

%======================================================================
\section*{Concluzii}
Am implementat semafoare numărătoare direct în kernelul xv6, ca trei apeluri de sistem noi. Am parcurs întregul drum al
unui apel de sistem: numărul în \texttt{syscall.h}, declarația din \texttt{user.h}, stub-ul generat de \texttt{usys.pl},
intrarea din tabelul \texttt{syscalls[]} și implementarea din \texttt{sem.c}. Un semafor de kernel se reduce la un
număr protejat de un spinlock și la mecanismul \texttt{sleep}/\texttt{wakeup}. Două reguli sunt esențiale: condiția se
verifică într-un \texttt{while}, pentru că \texttt{wakeup} trezește pe toți, iar spinlock-ul se eliberează înainte de a
dormi. Testele au confirmat ordonarea proceselor (copilul după părinte), alternarea strictă ping/pong și folosirea
semaforului ca lock (\texttt{mix lock}).

\end{document}
